\textbf{Part 1.}

\begin{description}
\item[1.1)] From the plot on the right we can get the equation of the
  straight line as
  \[ y = -0.76\cdot x + 7.32 \]
  \hfill[3 points]

  which taking to the power of 10 is equivalent to:
  $f_g = 10^{7.32}M_*^{-0.76}$, being $f_g$ the gas fraction. Multiplying
  by the stellar mass we finally get:
  \[ M_{gas} = 10^{7.32}\times M_*^{0.24} \]
  \[ a = 10^{7.32} = 2.09\times 10^{7}\ \text{[1 point]};\qquad
     b = 0.24\ \text{[1 point]} \]

\item[1.2)] \hfill[15 points]

  \begin{center}
  \begin{tabular}{cccccc}
  V & LogV & K & Zx & Zy & ZxZy\\
  79.4 & 1.8998 & $-16.8$ & $-1.1778$ & 1.3873 & $-1.6339$\\
  100.1 & 2.0004 & $-19.2$ & $-0.7831$ & 0.4970 & $-0.3892$\\
  158.5 & 2.2000 & $-21.3$ & $-0.0001$ & $-0.2819$ & 0.0000\\
  251.2 & 2.4000 & $-21.4$ & 0.7845 & $-0.3190$ & $-0.2502$\\
  316.2 & 2.5000 & $-24$ & 1.1765 & $-1.2834$ & $-1.5100$\\
  \hline
  Ave & 2.2001 & $-20.54$ & & sum ZxZy & $-3.7833$\\
  St. Dev & 0.2549 & 2.696 & & r & $-0.9458$
  \end{tabular}
  \end{center}

  The best-fit straight line is
  $K = -10\cdot\log_{10}(V_{max}) + 1.47$.
\end{description}

\textbf{Part 2.}

\begin{description}
\item[2.1)] \hfill[10 points]

  Galaxies are at the same distance, therefore the difference in apparent
  magnitudes is the same as in absolute magnitudes, and this relates to a
  luminosity ratio:
  \[ (k_1 - k_2) = (K_1 - K_2)
     = -2.5\times\log_{10}\!\left(\frac{L_1}{L_2}\right) \]
  \hfill[2 points]

  For having the same stellar populations, and ignoring interstellar
  extinction, the luminosity of each galaxy is proportional to its
  stellar mass with the same mass-to-light ratio for both galaxies, such
  that:
  \[ (k_1 - k_2) = -2.5\times\log_{10}
     \!\left(\frac{M_{*1}}{M_{*2}}\right) \]
  \hfill[2 points]

  The difference in magnitudes being 6, we obtain:
  \[ \left(\frac{M_{*1}}{M_{*2}}\right) = 10^{2.4} \]
  \hfill[2 points]

\item[2.2)] \hfill[4 points]

  Using the last result in combination to that of 1A:
  \[ \frac{M_{gas1}}{M_{gas2}}
     = \left(\frac{M_{*1}}{M_{*2}}\right)^{0.24} \]
  \hfill[2 points]
  \[ \frac{M_{gas1}}{M_{gas2}} = 10^{0.576} \]
  \hfill[2 points]

\item[2.3)] \hfill[6 points]

  We found the slope of the TF relation in 1B:
  $\frac{\Delta K}{\Delta\log(V_{max})} = 10$, with $\Delta K = 6$ for
  our galaxies: $\frac{V_{max1}}{V_{max2}} = 10^{0.6}$ [2 points]. Now
  consider the mass distribution as approximately
  spherically-symmetric, giving:
  \[ V_{max} = \sqrt{\frac{G\cdot M_{tot}(<R_{max})}{R_{max}}}, \]
  \hfill[2 points]

  which combined with the ratio just found two lines above leads to:
  \[ \frac{M_{tot1}}{M_{tot2}} = 10^{1.2} \]
  \hfill[2 points]
\end{description}

\textbf{Part 3.} \hfill[15 points]

Baryonic mass is
$M_{Baryonic} = \frac{4.39\times 10^{11}}{7.82}\,[M_\odot]$ and Dark
matter mass is given by
$M_{dm} = 6.82\times\frac{4.39\times 10^{11}}{7.82}\,[M_\odot]$. The
baryonic contribution needs to be decomposed into stellar and gaseous
masses, which is not doable analytically but must be faced numerically,
by tuning the numbers, knowing their sum and knowing the relation between
them from 1.1. The final table looks like this:

\begin{center}
\begin{tabular}{cccccc}
Galaxy & Apparent magnitude $k$ & $M_{gas}$ [$M_\odot$]
  & $M_{*}$ [$M_\odot$] & $M_{dm}$ [$M_\odot$] & $M_{tot}$ [$M_\odot$]\\
$G_1$ & 19.2 & $7.67\times 10^{9}$ & $4.85\times 10^{10}$
  & $3.83\times 10^{11}$ & $4.39\times 10^{11}$
\end{tabular}
\end{center}

Equation to be solved: $1 - k = 0.1415\,k^{0.24}$. This gives
$k = 0.8634$

\textbf{Part 4}

\begin{description}
\item[4.1)] \hfill[4 points]

  We basically found in 2.1 that
  $\left(\frac{M_{*1}}{M_{*2}}\right) = 10^{\frac{\Delta K}{-2.5}}$,
  being $\Delta K = -6$. In this new case, we have the extreme values of
  $\Delta K$ to be $-6.4$ and $-5.6$, therefore the extreme values for
  the exponent are these values divided by $-2.5$, then
  $e\in[2.24,\ 2.56]$

\item[4.2)] \hfill[10 points]

  We have to invert the equation from 1.2 to get:
  $\log_{10}(V_{max}) = -\frac{K}{10} + 0.1467$. Now we can get the
  difference between the measured values and linear-fit predictions,
  approximating numbers to the second digit:

  \begin{center}
  \begin{tabular}{ccccc}
  $K$ [mag] & $V_{max}$ [km/s] & $\log_{10}(V_{max})$ measured
    & $\log_{10}(V_{max})$ from linear fit
    & $|\Delta\log_{10}(V_{max})|$\\
  $-16.8$ & 79.4 & 1.9 & 1.83 & 0.07\\
  $-19.2$ & 100.1 & 2.0 & 2.07 & 0.07\\
  $-21.3$ & 158.5 & 2.2 & 2.28 & 0.08\\
  $-21.4$ & 251.2 & 2.4 & 2.29 & 0.11\\
  $-24.0$ & 316.2 & 2.5 & 2.55 & 0.05
  \end{tabular}
  \end{center}

  Taking two times the RMS of the residuals:
  $\sigma_{stat} = 0.157$

\item[4.3)] With the uncertainties, the expression
  $\log_{10}(V_{max}) = -\frac{K}{10} + 0.1467$ becomes:
  \[ \log_{10}(V_{max}) = -\frac{K}{10} + 0.1467\pm\frac{0.2}{10}
     \pm 0.157 \]
  \hfill[2 points]

  Writing the same expression down for galaxy 1 and 2, subtracting, and
  using $\Delta K = -6$ we get:
  \[ \log_{10}\!\left(\frac{V_{max1}}{V_{max2}}\right)
     = 0.6\pm 0.02\pm 0.157\pm 0.02\pm 0.157 \]
  \hfill[4 points]

  Which leads to the extreme values: $0.954$ and $0.246$.
  \hfill[2 points]

  Finally remembering that circular velocities scale as the square root
  of the total mass, we have to multiply this exponent by 2, getting the
  final answer:
  \[ g\in[0.49,\ 1.91] \]
  \hfill[2 points]
\end{description}
